\section{Limitations and Outlook}

The model is deliberately not a spacetime model. It has no metric, no causal horizon, no Bekenstein-Hawking area law, no semiclassical exterior region, no quantum extremal surface, no island, and no replica wormhole. It therefore does not address why the gravitational entropy prescription has the form it does. It only tests how much black-hole-like phenomenology can be reproduced by finite quantum statistical mechanics once a negative-heat-capacity density of states is supplied.

The core Hamiltonian is also engineered. The shell dimensions are chosen to realize a desired entropy profile, and the core Hamiltonian is block diagonal in those shells. This is enough for a control model, but it is not a natural analogue black hole.

The main open problem is therefore the next step:
\begin{equation}
  \text{find a natural finite Hamiltonian whose density of states supplies the convex window.}
\end{equation}
Possible directions include finite droplet models, attractive Bose-Hubbard clusters, long-range interacting spin systems, and finite systems with phase coexistence. In such systems a convex intruder can appear in the microcanonical entropy without imposing shell dimensions by hand.

A second open problem is to replace the engineered radiation-bin labels by a more physical bath or waveguide model. The present calculation shows why many outgoing channels matter: one-channel evaporation is rank-limited and creates dark subspaces. A more physical implementation would realize these channels as frequencies, angular modes, species labels, or spatial outgoing modes.

The correct conclusion is therefore modest. We have constructed a minimal geometry-free Hamiltonian evaporator that combines negative microcanonical heat capacity with unitary radiation emission in a controlled working window. This does not solve the black-hole information problem. It provides a comparison model for separating finite-dimensional quantum-statistical mechanisms from genuinely gravitational mechanisms.
